\documentclass[UTF8,11pt]{ctexart}
\usepackage[a4paper,margin=24mm]{geometry}
\usepackage{amsmath,amssymb,amsthm,mathtools,mathrsfs}
\usepackage[all]{xy}
\usepackage{xurl,hyperref,enumitem}
\usepackage{tikz}
\usetikzlibrary{arrows.meta,cd,decorations.markings,calc,patterns}
\hypersetup{hidelinks,breaklinks=true}
\setlength{\emergencystretch}{3em}
\allowdisplaybreaks
\newtheorem*{theorem}{Theorem}
\newtheorem*{thm}{Theorem}
\newtheorem*{lemma}{Lemma}
\newtheorem*{proposition}{Proposition}
\newtheorem*{prop}{Proposition}
\newtheorem*{corollary}{Corollary}
\newtheorem*{cor}{Corollary}
\newtheorem*{claim}{Claim}
\theoremstyle{definition}
\newtheorem*{definition}{Definition}
\newtheorem*{defn}{Definition}
\newtheorem*{remark}{Remark}
\newtheorem*{example}{Example}
\newcommand{\R}{\mathbb R}
\newcommand{\C}{\mathbb C}
\newcommand{\Z}{\mathbb Z}
\newcommand{\Q}{\mathbb Q}
\newcommand{\N}{\mathbb N}
\newcommand{\F}{\mathbb F}
\newcommand{\D}{\mathbb D}
\newcommand{\bB}{\mathbb B}
\newcommand{\bC}{\mathbb C}
\newcommand{\bR}{\mathbb R}
\newcommand{\bZ}{\mathbb Z}
\newcommand{\bN}{\mathbb N}
\newcommand{\bQ}{\mathbb Q}
\newcommand{\bD}{\mathbb D}
\newcommand{\bF}{\mathbb F}
\newcommand{\CP}{\mathbb{CP}}
\newcommand{\RP}{\mathbb{RP}}
\newcommand{\sO}{\mathscr O}
\newcommand{\sI}{\mathscr I}
\newcommand{\sF}{\mathscr F}
\newcommand{\sG}{\mathscr G}
\newcommand{\sH}{\mathscr H}
\newcommand{\sR}{\mathscr R}
\newcommand{\sM}{\mathscr M}
\newcommand{\sV}{\mathscr V}
\newcommand{\sU}{\mathscr U}
\DeclareMathOperator{\Hom}{Hom}
\DeclareMathOperator{\Ext}{Ext}
\DeclareMathOperator{\Tor}{Tor}
\DeclareMathOperator{\Spec}{Spec}
\DeclareMathOperator{\Proj}{Proj}
\DeclareMathOperator{\supp}{supp}
\DeclareMathOperator{\im}{im}
\DeclareMathOperator{\id}{id}
\DeclareMathOperator{\Tr}{Tr}
\DeclareMathOperator{\tr}{tr}
\DeclareMathOperator{\rank}{rank}
\DeclareMathOperator{\ord}{ord}
\DeclareMathOperator{\dist}{dist}
\DeclareMathOperator{\End}{End}
\DeclareMathOperator{\Aut}{Aut}
\DeclareMathOperator{\Gal}{Gal}
\DeclareMathOperator{\vol}{vol}
\DeclareMathOperator{\Cl}{Cl}
\DeclareMathOperator{\coker}{coker}
\DeclareMathOperator{\codim}{codim}
\DeclareMathOperator{\divergence}{div}
\DeclareMathOperator{\Ric}{Ric}
\DeclareMathOperator{\grad}{grad}
\DeclareMathOperator{\Hess}{Hess}
\newcommand{\abs}[1]{\left\lvert#1\right\rvert}
\newcommand{\sabs}[1]{\lvert#1\rvert}
\newcommand{\babs}[1]{\bigl\lvert#1\bigr\rvert}
\newcommand{\Babs}[1]{\Bigl\lvert#1\Bigr\rvert}
\newcommand{\bbabs}[1]{\biggl\lvert#1\biggr\rvert}
\newcommand{\BBabs}[1]{\Biggl\lvert#1\Biggr\rvert}
\newcommand{\norm}[1]{\left\lVert#1\right\rVert}
\newcommand{\snorm}[1]{\lVert#1\rVert}
\newcommand{\bnorm}[1]{\bigl\lVert#1\bigr\rVert}
\newcommand{\Bnorm}[1]{\Bigl\lVert#1\Bigr\rVert}
\newcommand{\bbnorm}[1]{\biggl\lVert#1\biggr\rVert}
\newcommand{\BBnorm}[1]{\Biggl\lVert#1\Biggr\rVert}
\newcommand{\widebar}[1]{\overline{#1}}
\newcommand{\nicefrac}[2]{\frac{#1}{#2}}
\newcommand{\linnprod}[2]{\langle#1,#2\rangle}
\newcommand{\myindex}[1]{#1}
\newcommand{\glsadd}[1]{}
\newcommand{\avoidbreak}{}
\newcommand{\sourceref}[1]{\textnormal{[原文：\nolinkurl{#1}]}}
\newcommand{\thmref}[1]{\sourceref{#1}}
\newcommand{\lemmaref}[1]{\sourceref{#1}}
\newcommand{\propref}[1]{\sourceref{#1}}
\newcommand{\corref}[1]{\sourceref{#1}}
\newcommand{\defnref}[1]{\sourceref{#1}}
\newcommand{\exampleref}[1]{\sourceref{#1}}
\newcommand{\exerciseref}[1]{\sourceref{#1}}
\newcommand{\figureref}[1]{\sourceref{#1}}
\newcommand{\sectionref}[1]{\sourceref{#1}}
\newcommand{\appendixref}[1]{\sourceref{#1}}
\newcommand{\stacksref}[2]{\href{https://stacks.math.columbia.edu/tag/#1}{#2 [#1]}}


\newcommand{\E}{\mathbb E}
\newcommand{\Prob}{\mathbb P}
\newcommand{\Reff}{\mathcal R_{\mathrm{eff}}}
\newcommand{\eps}{\varepsilon}
\DeclareMathOperator{\diam}{diam}
\DeclareMathOperator{\Var}{Var}
\DeclareMathOperator{\Crit}{Crit}
\newcommand{\dd}{\,\mathrm d}

\begin{document}
\section*{056\quad 正则图非平凡谱的Alon--Boppana型下界}
\subsection*{来源与计数目标}
MIT 18.217, \emph{Graph Theory and Additive Combinatorics}, Fall 2019, \S4.4, printed p.~88: ``Proof 2 (slightly weaker result) following Theorem 4.15. 课程讲义由学生据 Yufei Zhao 授课撰写，并由教师协助编辑。取得日期：2026-09-28。
\par\url{https://ocw.mit.edu/courses/18-217-graph-theory-and-additive-combinatorics-fall-2019/212e032b28e58a46e62f5a32d295821e_MIT18_217F19_full_notes.pdf}
\par\textbf{本文件只计一个问题：}固定 $d\geq3$，证明 $d$-正则图的非平凡特征值绝对值最大值至少为 $2\sqrt{d-1}-o(1)$；不是仅对 $\lambda_2$ 的更强版本。
\subsection*{作者命题与人类证明}
\paragraph{Setting (source \S4.4).}
\noindent\textit{以下背景与记号据所列原文章节摘编；不是逐字引文，也不是新增证明。}\par An $(n,d,\lambda)$-graph is a $d$-regular graph on $n$ vertices whose adjacency matrix has eigenvalues $d=\lambda_1\geq\cdots\geq\lambda_n$ and $\max\{|\lambda_2|,|\lambda_n|\}\leq\lambda$.
\begin{theorem}[Result proved in Proof 2 following Theorem 4.15]
\[\max\{|\lambda_2|,|\lambda_n|\}\geq2\sqrt{d-1}-o(1).\]
\end{theorem}
\begin{proof}[Proof 2 (slightly weaker result)]
We'll show that $\max\{|\lambda_2|,|\lambda_n|\}\geq2\sqrt{d-1}-o(1)$. This is an illustration of the trace method, also called the moment method. We have
\[\sum_{i=1}^n\lambda_i^{2k}=\operatorname{tr}(A^{2k}).\]
The right hand side is the number of closed walks of length $2k$ on $G$. Now, the number of closed walks of length $2k$ starting at a fixed vertex $v$ in a $d$-regular graph is at least the number of closed walks of length $2k$ starting at a fixed $v$ in an infinite $d$-regular tree. To see why this is true, note that given any walk on the infinite $d$-regular tree, we can walk in the same way on $G$ by assigning an orientation to each vertex. But $G$ may have more walks if it has cycles.

There are at least $C_k(d-1)^k$ closed walks of length $2k$ starting at a fixed $v$ in an infinite $d$-regular tree, where $C_k=\frac1{k+1}\binom{2k}{k}$ is the $k$-th Catalan number. Thus, the number of walks of length $2k$ on $G$ is at least $\frac{n}{k+1}\binom{2k}{k}(d-1)^k$. On the other hand,
\[d^{2k}+(n-1)\lambda^{2k}\geq\sum_{i=1}^n\lambda_i^{2k}.\]
Thus,
\[\lambda^{2k}\geq\frac1{k+1}\binom{2k}{k}(d-1)^k-\frac{d^{2k}}{n}.\]
The term $\frac1{k+1}\binom{2k}{k}$ is $(2-o(1))^{2k}$ as $k\to\infty$. Letting $k\to\infty$ and $k=o(\log n)$ as $n\to\infty$ gives us $\lambda\geq2\sqrt{d-1}-o(1)$.
\end{proof}
\subsection*{整理说明（非作者正文）}
从作者第二份证明原句提取其实际证明的结论，放入独立定理环境；上下文中的 $d$ 固定、$n\to\infty$ 及正则图条件已明确列出。未将第一份证明的较强题面配给第二份证明。允许引用邻接矩阵的谱分解、迹与闭游走计数、Catalan 数计数及其渐近性质；本题实际承担的是把树上的闭游走下界引入有限图，再选择增长阶数 $k=o(\log n)$ 分离平凡特征值。上述论证完整保留，树的例示图不是证明必要图，未转录。正文来自实际取得的 PDF 文本；未取得页面图像或下载原件。
\subsection*{可修改的简短标注（非作者正文）}
$c$：正则图的非平凡特征值下界。

$a$：邻接矩阵迹；闭游走；无限正则树；Catalan 数；高阶矩与渐近参数选择。

\texttt{cross\_theory\_status = yes}。将谱的偶次矩转成闭游走计数，用无限树上的组合路径给出统一下界，再从高阶矩返回有限图的特征值。位置：原文 Proof 2 的全部论证。
标签为非穷尽、可修改初稿。
\subsection*{许可与修改记录}
原作者及作品见来源栏；来源采用 CC BY-NC-SA 4.0。本题证摘录和附加整理沿用该许可。2026-09-28：增加题号、来源定位、独立排版、整理说明和初稿标注；数学内容的取材范围及排版变化见整理说明。
\end{document}
